\documentclass[11pt]{article}

\usepackage[utf8]{inputenc}
\usepackage[T1]{fontenc}
\usepackage{lmodern}
\usepackage[margin=1in]{geometry}
\usepackage{microtype}
\usepackage{parskip}
\usepackage[hidelinks]{hyperref}

\title{\vspace{-2em}\textbf{Teaching Statement}\\[0.3em]\large Alexandre Pavlov}
\author{}
\date{}

\begin{document}
\maketitle
\vspace{-3em}

\begin{center}
\small Université de Montréal \hspace{0.5em}\textbar\hspace{0.5em}
\href{mailto:alexandre.pavlov@umontreal.ca}{alexandre.pavlov@umontreal.ca} \hspace{0.5em}\textbar\hspace{0.5em}
\href{https://alexandrepavlov.com/}{alexandrepavlov.com}
\end{center}
\vspace{0.5em}

\textbf{Teaching philosophy.}
My approach to teaching is guided by two convictions. First, that economics is best learned through reasoning rather than memorization: a student who can derive the intuition behind an aggregate demand curve will outperform one who has memorized its shape. Second, that this kind of reasoning is built most reliably through dialogue rather than lecture --- that understanding emerges when students attempt to articulate a concept and have their reasoning gently corrected, refined, or extended. These convictions shape both the structure of my classes and the way I respond to students inside and outside them.

\textbf{Experience.}
During my PhD at the Université de Montréal, I have accumulated five years of teaching at both the undergraduate and graduate level. As a teaching assistant I have supported the first-year PhD macroeconomics course and the master's-level macroeconometrics course, work that has deepened my own command of advanced material; at the undergraduate level I have served as TA for applied econometrics and the third-year applied economics research workshop. Since 2022 I have also been the lecturer of record for first-year courses in macroeconomics and in economics and globalization, where I cover fundamental concepts including fiscal policy, monetary policy, and international trade. Across all courses I have taught, students have evaluated me with an average rating of 3.79 out of 4.00, against a departmental average of approximately 3.4, and my scores have improved year over year as I have iterated on my approach in response to student feedback.

\textbf{Building intuition before formalism.}
In macroeconomics in particular, students often arrive expecting to memorize a sequence of curves and equations. I structure lectures so that the formalism is always the destination, never the starting point. I will typically open a topic with a current policy question --- the design of a carbon price, the macroeconomic consequences of central-bank tightening, the trade-offs in fiscal stimulus during a recession --- and use it as the motivating problem against which we then develop the model. This gives the model a job to do: students see why we need it before we write it down, and they retain the mechanics better as a result. It also lets me draw naturally on my own research in climate macroeconomics and public finance, bringing live debates from the frontier of the discipline into the classroom in a way that undergraduates rarely see.

\textbf{Dialogue and the productive use of wrong answers.}
I encourage active participation by fostering a judgment-free environment where students feel comfortable speaking up. When a student offers an incorrect answer, I ask them to explain their reasoning and walk through it with them, rather than simply telling them they are wrong. The point is to make their thinking visible --- both to them and to the rest of the class --- so we can locate where the logic broke down. Over a semester this changes the room: students who were initially silent come to volunteer answers, because being wrong about a step is treated as part of learning rather than as failure. It also surfaces misunderstandings that a one-way lecture would have hidden until the exam.

\textbf{Adapting to diverse backgrounds.}
A first-year economics class at the Université de Montréal contains students with very different preparation. Some have completed advanced calculus, others have not seen a derivative in years; some are economics majors, others are taking the course as an elective from political science, history, or environmental studies. The student body is also linguistically diverse, with students whose first language ranges across French, English, and many others. I design my lectures and problem sets so that the core material is accessible to every student in the room, while leaving room for students who want to push further --- through optional reading suggestions, harder end-of-set questions, or deeper office-hour conversations. Having taught in both French and English, I am attentive to the way technical vocabulary travels (or fails to) across languages, and I make a habit of writing key terms on the board so that students following in their second language do not lose the thread.

\textbf{Feedback as an ongoing input.}
I treat student feedback as something I act on during the semester, not only after it. % [CUSTOMIZE: if you do mid-semester surveys, an open suggestion box, or any specific feedback ritual, name it here in one sentence — e.g., "I distribute a short anonymous mid-semester survey and discuss the responses with the class the following week."]
This has translated into concrete changes in pacing, problem-set design, and the balance between intuition-building and technical derivation, and is reflected in the year-over-year improvement in my evaluations.

% [CUSTOMIZE — STRONGLY RECOMMENDED: insert one short anecdote here (3–5 sentences) of a specific moment that exemplifies your philosophy. The single most memorable element of any teaching statement is a concrete story. For example: a student who arrived believing they "weren't a math person" and who, by working through a derivation aloud in office hours, discovered they could in fact reason their way through it; a class where a misconception about monetary policy surfaced during discussion and produced a richer lecture than the one you had planned; a student who told you (or wrote in evaluations) that your course changed how they thought about economics. Choose the story that you would most want a hiring committee to remember.]

\textbf{Looking ahead.}
Looking forward, I am especially eager to develop courses that bring my research interests into the classroom. Climate macroeconomics is one of the most policy-relevant frontiers of the discipline today, and few undergraduates encounter it formally; I would welcome the opportunity to design a course introducing students to the social cost of carbon, optimal climate policy, and the macroeconomic implications of the energy transition, drawing on both the canonical literature and live debates. At the graduate level, I am well-positioned to teach core macroeconomics, public finance, and applied climate-economics modules, including the computational tools (Julia, Python, MATLAB, Fortran) that increasingly define applied work in the field. % [CUSTOMIZE: if there are specific courses you would most like to teach, or course numbers/areas at the institution you are applying to, name them here.]
More broadly, I see teaching and research as continuous with one another: the questions that animate my research --- how policy should respond to climate risk, how heterogeneity shapes optimal taxation, how rules and discretion interact in long-horizon decisions --- are precisely the questions I want my students to be able to reason about by the end of a semester with me.

\end{document}
